The method is motivated by a state-observability problem. An incorrect program may follow from an inadequate task interpretation, an incorrect implementation of an adequate interpretation, or both. Direct NL-to-Code exposes only the final program and cannot show where the failure entered. We therefore place an explicit semantic state between problem $x$ and program $C$, revise that state before implementation, and hold it fixed during program repair.

\subsection{Overview}

Figure~\ref{fig:framework} divides generation into two connected stages. SAL induces an initial specification $s^{(0)}$. At semantic iteration $t$, it proposes a field-level patch $\Delta^{(t)}$, merges the patch with the current specification to obtain $\hat{s}^{(t+1)}$, and accepts the candidate only if it satisfies both safety and improvement conditions. An accepted candidate directly replaces the current state; a rejected candidate leaves $s^{(t)}$ unchanged and stops SAL. The current specification at termination is frozen as semantic contract $s^\star$, from which executable property set $P^\star$ is compiled.

IRL then generates an initial candidate set $\mathcal C^{(0)}$ from $(x,s^\star)$ and uses public feedback to select the unique initial state $C^{(0)}$. At round $k$, only current program $C^{(k)}$ is verified and repaired. The resulting repair candidates form the round-specific set $\mathcal C^{(k+1)}$, from which the system selects $C^{(k+1)}$. Unselected candidates from earlier rounds do not carry over. The complete process records semantic states, implementation states, and their evidence in trace $\tau$:
\begin{equation}
\begin{aligned}
x &\rightarrow s^{(0)}\rightarrow(q^{(0)},\mathrm{SAS}^{(0)})\rightarrow\cdots
\rightarrow(\Delta^{(t)},\hat{s}^{(t+1)})\\
&\rightarrow \mathrm{Accept}^{(t+1)}\rightarrow s^{(t+1)}\rightarrow\cdots
\rightarrow s^\star\rightarrow P^\star,\\
(x,s^\star)&\rightarrow\mathcal C^{(0)}\rightarrow C^{(0)}\rightarrow\cdots\rightarrow C^{(k)}\\
&\rightarrow(F_v^{(k)},F_p^{(k)})\rightarrow\mathcal C^{(k+1)}
\rightarrow C^{(k+1)}\rightarrow\cdots\rightarrow\tau.
\end{aligned}
\label{eq:transition}
\end{equation}
This boundary does not treat $s^\star$ as a proof of task intent. The contract is an LLM-induced, inspectable target; SAS is a directional evaluation signal; and only a limited subset of obligations can be compiled into executable clauses. Algorithm~\ref{alg:dual_loop} gives the complete control flow.

\begin{algorithm}[t]
\caption{Semantic-contract-based dual-loop generation}
\label{alg:dual_loop}
\begin{algorithmic}[1]
\REQUIRE problem $x$; SAL budget $T_s$; IRL budget $T_r$
\ENSURE final program $C$ and trace $\tau$
\STATE \textbf{Stage 1: semantic alignment}
\STATE $s^{(0)}\gets\algcall{InduceSpec}{x}$
\STATE $(q^{(0)},\mathrm{SAS}^{(0)})\gets\algcall{Judge_{\mathbf w}}{x,s^{(0)}}$
\FOR{$t=0,\ldots,T_s-1$}
  \IF{$\neg\algcall{ShouldRefine}{s^{(t)},q^{(t)},\mathrm{SAS}^{(t)},r}$}
    \STATE \textbf{break}
  \ENDIF
  \STATE $\Delta^{(t)}\gets\algcall{ProposePatch}{x,s^{(t)},q^{(t)}}$
  \STATE $\hat{s}^{(t+1)}\gets\algcall{Merge}{s^{(t)},\Delta^{(t)}}$; judge candidate
  \IF{$\mathrm{Safe}^{(t+1)}\land\mathrm{Improve}^{(t+1)}$}
    \STATE $s^{(t+1)}\gets\hat{s}^{(t+1)}$
  \ELSE
    \STATE $s^{(t+1)}\gets s^{(t)}$; \textbf{break}
  \ENDIF
\ENDFOR
\STATE $s^\star\gets$ current specification at SAL termination; freeze $s^\star$
\STATE $P^\star\gets\algcall{CompileClauses}{s^\star}$; initialize trace $\tau$
\STATE \textbf{Stage 2: implementation and repair}
\STATE $\mathcal C^{(0)}\gets\algcall{GenerateCandidates}{x,s^\star,N_0}$
\STATE $C^{(0)}\gets\algcall{SelectByFeedback}{\mathcal C^{(0)},x,s^\star,P^\star}$
\FOR{$k=0,\ldots,T_r$}
  \STATE $F_v^{(k)}\gets\algcall{Verify_{fb}}{x,C^{(k)}}$
  \IF{$\algcall{Pass}{F_v^{(k)}}\lor k=T_r$}
    \STATE \textbf{return} $(C^{(k)},\tau)$
  \ENDIF
  \STATE $F_p^{(k)}\gets\algcall{EvaluateClauses}{P^\star,x,C^{(k)},F_v^{(k)}}$
  \FOR{$i=1,\ldots,N_r$}
    \STATE $C_i'^{(k+1)}\gets\algcall{Repair_i}{x,s^\star,C^{(k)},F_v^{(k)},F_p^{(k)}}$
  \ENDFOR
  \STATE $\mathcal C^{(k+1)}\gets\{C_i'^{(k+1)}\mid 1\leq i\leq N_r\}$
  \STATE $C^{(k+1)}\gets\algcall{SelectByFeedback}{\mathcal C^{(k+1)},x,s^\star,P^\star}$
  \STATE retain $C^{(k)}$ if the selected result is invalid or unchanged; update $\tau$
\ENDFOR
\end{algorithmic}
\end{algorithm}

\subsection{Semantic Alignment Loop}

SAL contains a specification inducer, a prompt-driven semantic judge, a field-level revision module, and an acceptance gate. The six central field groups correspond to recurrent failure surfaces identified in prior defect and specification-misunderstanding studies \citep{ref_bug_study,ref_codegen_error_taxonomy,ref_mufix,ref_pseudoeval}: task summary $s_{\mathrm{sum}}$, I/O protocol $s_{\mathrm{io}}$, constraints $s_{\mathrm{cons}}$, rules $s_{\mathrm{rules}}$, edge cases $s_{\mathrm{edge}}$, and candidate properties $s_{\mathrm{prop}}$:
\begin{equation}
s\triangleq(s_{\mathrm{sum}},s_{\mathrm{io}},s_{\mathrm{cons}},s_{\mathrm{rules}},s_{\mathrm{edge}},s_{\mathrm{prop}}).
\label{eq:spec}
\end{equation}
The decomposition lets the system localize a protocol error, missing constraint, unsupported rule, or boundary omission without regenerating the whole interpretation. Candidate properties are obligation statements, not trusted assertions; only clauses matching supported and reliably evaluable schemas can later become executable feedback.

The semantic judge uses the same backbone at temperature zero with a dedicated rubric and JSON output schema. Given $x$ and $s$, it returns coverage, faithfulness, and precision profile $q(x,s)$ together with aggregate SAS:
\begin{equation}
\begin{aligned}
(q(x,s),\mathrm{SAS}(x,s))&=\mathrm{Judge}_{\mathbf w}(x,s),\\
\mathbf w&=(0.4,0.4,0.2).
\end{aligned}
\label{eq:sas}
\end{equation}
Coverage measures whether important obligations and edge cases are retained; faithfulness penalizes unsupported assumptions; and precision measures whether obligations are operational enough to guide generation and repair. Outputs are integers in $[0,100]$. The pipeline validates and records the aggregate returned under the rubric rather than recomputing a post-hoc dot product, and it does not treat SAS as an independent ground-truth label. Coverage and faithfulness receive equal primary weight because either omission or invention can redirect generation; precision receives lower weight because operational detail cannot compensate for missing or distorted semantics. Section~\ref{app:sensitivity} reports the sensitivity experiment supporting this fixed setting.

The judge additionally returns grounded issue evidence, issue types, at most three target fields, protected fields, and a field-level edit plan. A patch may change only target fields, may not change protected fields, and must produce a material semantic difference. Let $\mathcal F$ denote the field set and $\mathcal T^{(t)}$, $\mathcal G^{(t)}$, and $\Pi^{(t)}$ denote the target set, protected set, and edit plan. A candidate outside these constraints is rejected before re-evaluation.

The acceptance gate separates non-regression from grounded progress:
\begin{equation}
\mathrm{Accept}^{(t+1)}=1\iff
\mathrm{Safe}^{(t+1)}\land\mathrm{Improve}^{(t+1)}.
\label{eq:accept}
\end{equation}
$\mathrm{Safe}$ requires a valid material patch, non-decreasing SAS and faithfulness, no increase in unsupported issues, and bounded degradation in missing-issue and confidence signals. $\mathrm{Improve}$ requires at least one grounded gain, such as an SAS increase of at least $\delta$, reduced issue pressure, improved precision without coverage loss, or fewer missing or blocking obligations. Appendix~\ref{app:gate} gives the complete predicates. If Eq.~\eqref{eq:accept} holds, $s^{(t+1)}=\hat{s}^{(t+1)}$; otherwise $s^{(t+1)}=s^{(t)}$ and SAL stops. The main configuration uses $\delta=1$. Refinement is attempted only when $\mathrm{SAS}^{(t)}<r$, the judge reports a grounded material issue, and a field-specific plan exists. We use $r=90$, at most three attempts, and stop after one rejected or skipped revision.

No additional specification candidates are generated or ranked after this trajectory. The current specification at SAL termination is directly frozen as $s^\star$. A conservative compiler then extracts the executable subset
\begin{equation}
P^\star=\bigcup_{\ell\in\mathrm{Fields}(s^\star)}\mathrm{Compile}(\ell).
\label{eq:compile}
\end{equation}
Supported schemas cover ordering, multiset preservation, output length, Boolean protocols, numeric and non-negative output, and modulo ranges. General optimality and algorithmic-correctness claims remain textual. A recognized clause may still abstain if observed and expected values do not support a reliable check.

\subsection{Implementation Repair Loop}

The contract-conditioned generator receives both $x$ and $s^\star$, produces at most $N_0$ initial candidates, and forms $\mathcal C^{(0)}$. Each candidate is evaluated on public feedback tests. A fixed lexicographic rank considers complete public pass, pass ratio and count, failure progress, error type, weighted property violations, empty output, violated specification items, output-length difference, and matching lines. Private tests never participate. The selected program $C^{(0)}$ is the unique initial IRL state; unselected initial candidates are not repaired later.

At round $k$, the feedback verifier executes current program $C^{(k)}$ on public tests and returns structured evidence $F_v^{(k)}$, including pass status, error type, failing input, observed output, available expected output, and partial progress. SAS evaluates semantic state before implementation; $F_v^{(k)}$ evaluates executable behavior. SAL, initial selection, and IRL never query private tests. The final evaluator uses the complete public-and-private suite only after a single final program has been selected.

Let $P^\star=\{\phi_i\}_{i=1}^{N_p}$, where $N_p$ is the number of successfully compiled clauses that can be evaluated reliably. The property oracle evaluates those clauses using the problem, current program, and verifier evidence:
\begin{equation}
\begin{aligned}
F_v^{(k)}&=\mathrm{Verify}_{fb}(x,C^{(k)}),\\
F_p^{(k)}&=\mathrm{EvaluateClauses}(P^\star,x,C^{(k)},F_v^{(k)}).
\end{aligned}
\label{eq:feedback}
\end{equation}
$F_v^{(k)}$ reports concrete execution behavior, whereas $F_p^{(k)}$ records violated executable obligations and their evidence. Satisfied clauses and clauses that cannot be judged reliably produce no violation feedback. Neither feedback source may modify $s^\star$.

The repair module then generates a round-specific candidate set from the unique current program:
\begin{equation}
\begin{aligned}
C_i'^{(k+1)}&=\mathrm{Repair}_i(x,s^\star,C^{(k)},F_v^{(k)},F_p^{(k)}),quad 1\leq i\leq N_r,\\
\mathcal C^{(k+1)}&=\{C_i'^{(k+1)}\mid 1\leq i\leq N_r\},\\
C^{(k+1)}&=\mathrm{SelectByFeedback}(\mathcal C^{(k+1)},x,s^\star,P^\star).
\end{aligned}
\label{eq:repair}
\end{equation}
Selection uses the same public-feedback ranking as initial selection. Earlier candidate sets and their unselected members do not accumulate. If the selected result is unparsable or identical to $C^{(k)}$, no program-state change is recorded. IRL stops when the current program passes public feedback or the repair budget is exhausted:
\begin{equation}
\mathrm{Stop}_{\mathrm{IRL}}^{(k)}=1\iff
\mathrm{Pass}(F_v^{(k)})\lor(k=T_r),\qquad 0\leq k\leq T_r.
\label{eq:stop}
\end{equation}
The final program is then evaluated on held-out private tests. Trace labels summarize visible semantic-side, implementation-side, or mixed evidence; they support inspection but are not causal diagnoses, and the conservative policy returns \texttt{unknown} when evidence remains mixed.
